% Verbatim copy of /Users/alingalatan/workspaces/grothe/handoff/manuscript/legacy/ch25-parts/30-response-feasibility.tex (file dated 2026-10-25). Only this line was added.
\subsection{Keeping the fixed first moment: remote feasibility and a four-block response dual}
\label{sec:fixed-first-remote-feasibility}
\paragraph{Provenance (IDE note, 20 October).} Lower-bound GPT response-feasibility package (\code{conjK-response-feasibility-20261020.zip}, archived in
\pathref{gpt-reviews/response-feasibility-20261020/}). IDE replay \code{python3 -I verify\_all.py} from an unrelated directory: all exact algebra audits PASS (1.5\,s).
IDE check \st{P} of the cut: $\langle q,k\rangle=p\langle q,F\rangle=pe$ because $q=\E[F\mid X]$ is the orthogonal remote projection and $q\in\operatorname{ran}Q$;
$\langle Z,h\rangle=t$ from $Ph=tZ$; then $|f_\pm|\le1$ and conditioning on $X$. The energy ceiling $e\le4\varepsilon^2B(A-t)$ is the optimum $\lambda=p/(2B)$ of the
tangent form, giving $0.04626$ at $t=0.84$. The package's smoothed cap $858733/625000\approx S+1.93\cdot10^{-5}$ matches the corrected budget of \autoref{ch:K-state}.

\paragraph{Status and scope.}
The analytic statements below are proved [P]; their rational constants and
finite algebra audits have a portable exact replay [C]. They constrain the
\emph{original response problem}, with $Ph=tZ$ and $F=\varepsilon Qk$.
They do not assert a new inequality for every raw-K cubic, prove K, or
give a new Grothendieck endpoint. All shared sectors may be present.

Write $P=P_{1,0}$, $Q=P_{2,1}$, $p=1/\varepsilon$ and
$h=(f_++f_-)/2$, $k=(f_+-f_-)/2$, where $|f_\pm|\le1$.
Assume $Ph=tZ$, $t\ge0$, and $F=\varepsilon Qk$.
Let $q=\mathbb E[F\mid X]$ be the remote cubic, with $Z$ independent of
the remote Gaussian vector $X$, and put $e=\|q\|_2^2$.
Every other sector of $F$ is unrestricted. Scalar products are
conjugate-linear in the first argument.

\subsubsection{An exact cut on the remote law [P]}
Define the unit-variance Rice mean
\[
 \rho(u)=\mathbb E_Z|Z+\sqrt u|,\qquad u\ge0.
\]
For every $\lambda\ge0$,
\begin{equation}\label{eq:remote-feasibility-cut}
 \boxed{\quad t+\lambda pe\le
       \mathbb E_X\rho(\lambda^2|q|^2).\quad}
\end{equation}
In particular, this is a linear constraint on the law of $|q|^2$:
\[
 \mathbb E\bigl[\rho(\lambda^2|q|^2)-\lambda p|q|^2\bigr]\ge t.
\]
It requires no independence of $q$ from the remote drivers of the other
sectors, and no fourth- or sixth-moment assumption.

\begin{proof}
Since $q\in\operatorname{ran}Q$ and is the orthogonal remote component of
$F$, $\langle q,k\rangle=p\langle q,F\rangle=pe$.
Also $\langle Z,h\rangle=t$. Consequently,
\begin{align*}
 t+\lambda pe
 &=\frac12\operatorname{Re}\mathbb E\bigl[
    (\bar Z+\lambda\bar q)f_+
   +(\bar Z-\lambda\bar q)f_-\bigr]\\
 &\le\frac12\mathbb E\bigl[|Z+\lambda q|+|Z-\lambda q|\bigr]
 =\mathbb E_X\rho(\lambda^2|q|^2).
\end{align*}
The last equality conditions on $X$ and uses the circular Gaussian law
of $Z$. Circularity of $q$ itself is unnecessary.
\end{proof}

More generally, any remote test $v\in\operatorname{ran}Q$ gives
$t+p\operatorname{Re}\langle v,q\rangle\le\mathbb E\rho(|v|^2)$.
Thus a negative support gap is a certificate that the proposed response
moments cannot occur, not a large negative value of the original norm.

The derivative representation
\begin{equation}\label{eq:unit-rice-derivative}
 \rho'(u)=\frac1{2\sqrt\pi}\int_0^1
       e^{-us}\sqrt{\frac{1-s}{s}}\,ds
\end{equation}
shows that $\rho$ is increasing and concave. The formula follows from
Gaussian convolution of the norm (equivalently from its Laplace integral);
differentiation is justified by the integrable beta weight.
In particular $\rho(0)=\sqrt\pi/2$, $\rho'(0)=\sqrt\pi/4$.
Jensen also gives the exact necessary envelope
\[
 t\le\inf_{\lambda\ge0}
       \{\rho(\lambda^2e)-\lambda pe\}.
\]
Unlike Gaussian domination, this compares with a \emph{constant-modulus}
remote law and is a direct consequence of concavity.

\subsubsection{A global rational tangent and explicit energy ceilings [P,C]}
Set
\[
 A=\frac{886762780}{10^9},\qquad B=\frac{432306903}{10^9}.
\]
The portable rational checker proves
\begin{equation}\label{eq:remote-rice-tangent}
 \rho(u)\le A+Bu\qquad(u\ge0).
\end{equation}
Substituting $\lambda=p/(2B)$ into
\eqref{eq:remote-feasibility-cut} gives the dimension-free bound
\begin{equation}\label{eq:remote-energy-ceiling}
 \boxed{\quad e\le4\varepsilon^2B(A-t).\quad}
\end{equation}
At $\varepsilon=1891/2500$, exact rational comparisons give
\[
\begin{array}{c|c}
\text{condition}&\text{necessary remote energy ceiling}\\\hline
t\ge0.840&e<0.046266\\
t\ge0.845&e<0.041319\\
t\ge0.851&e<0.035383
\end{array}
\]
These restrictions hold with all simultaneous $A/B$, even and diagonal
sectors. They concern $e=\|q\|^2$, not the full energy $\|F\|^2$.
The separate elementary bound
$\|F\|^2\le\varepsilon^2(1-t^2)$ remains valid; its maximum on the
central interval is $0.168438298624$. \emph{Superseded as a domain cap} by the all-sector energy ellipsoid
\eqref{eq:sector-ellipsoid} of \S\ref{sec:sector-rice-feasibility}, which gives
$\|F\|^2<0.081884$ for $t\ge0.84$ and $\|F\|^2<0.071981$ for $t\ge0.8456$ (all sectors,
every dimension); keep the ellipsoid itself where possible, since its sector weights are
stronger than the total cap.

\paragraph{Exact tangent certificate.}
Let $R=\sqrt\pi/2$ and $u_0=1/10$. From
\eqref{eq:unit-rice-derivative},
\[
 \frac{\rho'(u_0)}R
 =\frac12\sum_{n\ge0}\frac{(-u_0)^n(1/2)_n}{n!(2)_n},
 \qquad
 \frac{\rho(u_0)}R
 =1+\frac12\sum_{n\ge0}
       \frac{(-1)^n u_0^{n+1}(1/2)_n}{(n+1)n!(2)_n}.
\]
The terms decrease in modulus, so consecutive truncations bracket both
quantities. The checker uses 15 and 16 terms, bounds $R$ by rational
squares, and bounds $\pi$ using alternating series in
$\pi=16\arctan(1/5)-4\arctan(1/239)$.
It verifies $B>\rho'(u_0)$ and
$A>\rho(u_0)-u_0\rho'(u_0)$. Concavity then proves
\eqref{eq:remote-rice-tangent} on the entire unbounded domain.
There is no radial cutoff or numerical quadrature in this certificate.

\paragraph{A polynomial moment version.}
Since $\rho^{(4)}\le0$, its cubic Taylor polynomial at zero is a global
majorant. Thus, with $M_j=\mathbb E|q|^j$,
\begin{equation}\label{eq:remote-moment-feasibility}
 t+\lambda pe\le\frac{\sqrt\pi}{2}
 \left(1+\frac{\lambda^2e}{2}
       -\frac{\lambda^4M_4}{16}
       +\frac{\lambda^6M_6}{96}\right).
\end{equation}
This is another necessary condition, with no Taylor truncation error of
unknown sign. The full Rice cut is preferable when the sixth moment makes
this polynomial loose. In a law LP, an upper enclosure of the Rice
function yields a safe necessary relaxation; a lower enclosure cannot
be substituted on the right to exclude a candidate.

\paragraph{A reusable exact kernel table.}
The companion module constructs 227 rational tangent planes and takes
their minimum, giving a global upper envelope of $\rho$ on $[0,\infty)$.
On $[0,1]$ its error is below $9\cdot10^{-7}$. Indeed
$|\rho''|\le\sqrt\pi/16$, the nodes there have spacing $1/128$, and
the exact rounding and series error is added to the nearest-tangent
error. Every plane is global, so the far tail needs no extrapolation
assumption. For an LP in the normalized variable $u=|q|^2/e$, the
correct coefficient is $\overline\rho(\lambda^2 e u)-\lambda peu$.

\subsubsection{An excluded coupled--diffuse family, robust under dependent perturbations [P,C]}
Let $X,Y,W_1,\ldots,W_N$ be independent standard proper Gaussians, remote
from $Z$, and put
\[
 q_N=\frac7{100}X\bigl[(|X|^2-2)+i(|Y|^2-1)\bigr]
       +\frac4{25\sqrt N}\sum_{j=1}^N H(W_j).
\]
For every integer $N\ge20$, every remote $\Delta\in H_{2,1}$ with
$\|\Delta\|_2\le1/1000$, and every $t\ge1057/1250$, there are
\emph{no} responses with
\[
 |f_\pm|\le1,\qquad Ph=tZ,\qquad
 \mathbb E[\varepsilon Qk\mid X,Y,W_1,\ldots,W_N,\text{other remote coordinates}]
       =q_N+\Delta.
\]
The perturbation may depend arbitrarily on the displayed remote
variables and additional remote Gaussians. Every visible sector remains
unrestricted. This is an exclusion from the original norm problem,
not a claim that the corresponding raw-K payoff is negative.

\begin{proof}
Set $c=7/100$, $g=4/25$, $x=1/N$, and $e=3c^2+g^2=403/10000$.
Exponential moments of the coupled core are $3c^2$, $210c^4$,
$98166c^6$. Independence and circularity of the two summands give
\begin{align*}
 M_4(x)&=210c^4+(2+42x)g^4+12c^2g^2,\\
 M_6(x)&=98166c^6+(6+378x+10968x^2)g^6\\
 &\hspace{12mm}+27c^2(2+42x)g^4+1890g^2c^4.
\end{align*}
In particular $4<M_4/e^2<6$ throughout $0\le x\le1/20$.

The cubic Hermite interpolant of $\rho$ with double nodes at $0$ and $4$
is a global majorant: the remainder is
$\rho^{(4)}(\xi)u^2(u-4)^2/4!\le0$.
Outward rational coefficients give the simpler majorant
\[
 \rho(u)\le P_3(u):=
 0.886226925453(1+u/2)-0.046344069u^2+0.003327750036u^3
 \qquad(u\ge0).
\]
All displayed decimals are exact rationals, certified by the same
alternating series and rational square bounds as above.
For $\lambda=3/2$, the checker verifies
\[
 \mathbb E P_3(\lambda^2|q_N|^2)
   -\frac{1057}{1250}-\lambda pe<-\frac9{20000}.
\]
The left side's majorant is a convex quadratic in $x=1/N$; checking
$x=0,1/20$ proves this for every $N\ge20$, including the diffuse limit.
Apply the general remote-test inequality with $v=\lambda q_N$.
Replacing $q_N$ by $q_N+\Delta$ on its left costs at most
\[
 \lambda p\sqrt e\,\|\Delta\|_2
 <\frac32\frac{2500}{1891}\frac{201}{1000}\frac1{1000}
 <\frac1{2500}.
\]
The support gap is still less than $-1/20000$, a contradiction.
\end{proof}

\subsubsection{Retaining the full first-chaos constraint in the response dual [P]}
Put $B_0=\beta Z(|Z|^2-2)$, $L=\ell Z+B_0$, $m=tZ$.
Let $R_t(F)$ be the supremum of the response-dependent terms after the
row constant, retaining both $Ph=m$ and $Qk=pF$ and both negative
squares. For $u,d\in\operatorname{ran}P$ and
$v,C\in\operatorname{ran}Q$, completion of squares gives
\begin{align}
R_t(F)\le D_{t,F}(u,v,C,d):={}&
\mathbb E|L+F-u-v-C-d|\nonumber\\
&+\mathbb E|L-F+u+v-C-d|
 +\|u\|^2+p\|C\|^2\nonumber\\
&-2p\operatorname{Re}\langle C,B_0\rangle
 +2p\operatorname{Re}\langle v,F\rangle
 +2\operatorname{Re}\langle d,m\rangle-p\|F\|^2.
\label{eq:fixed-first-response-dual}
\end{align}
No stationarity is used. The common first-chaos multiplier $d$ is free
because $Ph=m$ is an equality constraint.

\begin{proof}
Use
$-\|Pk\|^2\le\|u\|^2-2\operatorname{Re}\langle u,k\rangle$
and
\[
-\varepsilon\|Qh+pB_0\|^2
\le p\|C\|^2-2\operatorname{Re}\langle C,h\rangle
                 -2p\operatorname{Re}\langle C,B_0\rangle.
\]
Add the two zero terms
$2\operatorname{Re}\langle v,pF-Qk\rangle$ and
$2\operatorname{Re}\langle d,m-Ph\rangle$.
The coefficients of $f_+$ and $f_-$ are exactly the two expressions
inside the norms in \eqref{eq:fixed-first-response-dual}.
\end{proof}

The multiplier $d=sZ$ alone is equivalent to choosing a different visible
linear profile for this particular $F$; it is not a new price-free
improvement to a profile already optimized for every $F$ separately.
Its usefulness is the adaptive order: the trial can be chosen for the
joint configuration \emph{before} taking a supremum over that configuration.
Remote components of $d$ vanish in the pure odd baseline by parity, but
can interact with simultaneous even sectors.

\paragraph{Clipped Gram form.}
Let $w_\pm=L\pm F$, $\rho_\pm=\max(|w_\pm|,\tau)$, and define
$k_\tau=(w_+/\rho_+-w_-/\rho_-)/2$,
$h_\tau=(w_+/\rho_++w_-/\rho_-)/2$.
The four residual blocks are
\[
 r=(Pk_\tau,\ Qk_\tau-pF,\ Qh_\tau+pB_0,\ Ph_\tau-m).
\]
The quadratic form on the trial variables is
\[
 H(u,v,C,d)=\|u\|^2+p\|C\|^2+
 \frac12\mathbb E\left[
 \frac{|u+v+C+d|^2}{\rho_+}
 +\frac{|u+v-C-d|^2}{\rho_-}\right].
\]
With the usual
$B_\tau=\mathbb E[(\tau-|w_+|)_+^2+(\tau-|w_-|)_+^2]/(2\tau)$,
\[
 R_t(F)\le\Lambda(F)+B_\tau-r^*H^{-1}r
\]
on finite linearly independent trial spaces. Rational trials replace the
inverse by $2\operatorname{Re}(x^*r)-x^*Hx$.

Partition the matrix into the old three blocks and the new first-moment
block: $H=\left(\begin{smallmatrix}H_0&E\\E^*&J\end{smallmatrix}\right)$,
$r=(r_0,s)$. The additional optimized deduction is exactly
\[
 (s-E^*H_0^{-1}r_0)^*
 (J-E^*H_0^{-1}E)^{-1}
 (s-E^*H_0^{-1}r_0)\ge0.
\]
It is not the independently computed square of $Ph_\tau-m$; the
cross block is necessary to avoid counting the same adjustment twice.

\subsubsection{Implementation and remaining obligation}
The new necessary cuts can shrink the law relaxation only in a row that
retains the original response meanings of $t$ and $q$. A raw-K supremum
has no such feasibility assumption. Inserting these cuts into it would
be invalid. Equally, the $q$ in a priced auxiliary potential must first be
identified with this remote orthogonal projection, with its correct scale.

The accompanying diagnostic uses the full eighteen-dimensional complex
$H_{2,1}$ core on $(Z,X_1,X_2)$ plus an independent grade-three Gaussian
limit variable. It compares the three-block and four-block duals and
integrates the Gaussian buffer analytically. This is a limit-model test,
not an assertion that arbitrary finite cubics admit a buffer.
All diagnostic values are floats [F]. No unresolved all-tensor cap or
complete row cover is certified by them.

The portable checker verifies the rational tangent, three energy ceilings,
300 exact response square identities, 100 Schur identities, 324 Hermite
covariance/pseudocovariance identities, and the endpoint budget arithmetic.
The general proofs are the arguments above, not extrapolations from those
finite tests.

\subsection{Response-feasible remote laws: exact scalar reduction and a joint projection constraint}
\paragraph{Provenance (IDE note, 20 October).} Lower-bound GPT feasible-laws package (\code{conjK-feasible-laws-and-projections-20261020.zip}, archived in
\pathref{gpt-reviews/feasible-laws-20261020/}). IDE replay \code{python -I verify\_all.py} in an isolated venv with the pinned \code{python-flint 0.9.0}, from an unrelated
directory: ALL CLAIMS PASS (fresh integrals and tails, 88\,s). IDE check \st{P} of \eqref{eq:flr-response-stability}: the support bound holds for any remote $v$, since
$\langle v,k\rangle=\langle v,\bar k\rangle$; the Hessian is at most $2\rho'(0)=\sqrt\pi/2$ by concavity of $\rho$; then the square completion at $v=\lambda q$. In the
two-response form the $u$-gradients of $\rho(|u\pm v|^2)$ cancel at $u=0$. The three remote caps concern the remote component only.

\label{sec:response-feasible-law-reduction}
\paragraph{Scope.}
Analytic claims are labelled [P], directed computations [C], and exploratory
values [F]. No new Grothendieck endpoint is claimed. Set
\[
 \varepsilon=1891/2500,\quad p=\varepsilon^{-1},\quad
 L=Z(699/1250+(179/1250)|Z|^2),\quad
 f(u)=2\mathbb E|L+\sqrt u|,\quad \rho(u)=\mathbb E|Z+\sqrt u|.
\]
Here $Z$ is a standard proper complex Gaussian, independent of the remote
Gaussian vector $X$. For the actual bounded responses write
$h=(f_++f_-)/2$, $k=(f_+-f_-)/2$, $Ph=tZ$,
$F=\varepsilon Qk$, and $q=\mathbb E[F\mid X]$.
The cut already proved in the preceding response-feasibility section is
\begin{equation}\label{eq:flr-cut}
 t+\lambda p e\le\mathbb E\rho(\lambda^2|q|^2),\qquad
 e=\|q\|_2^2,\quad \lambda\ge0.
\end{equation}
All visible sectors of $F$ remain unrestricted. The bounds below on
$\mathbb E f(|q|^2)-pe$ concern the remote contribution; they cannot
replace a bound on the full $\Lambda(F)$.

\subsubsection{The scalar law problem collapses to two points [P]}
Let $m>1$, $M>m^2$. Put
\[
 A=\frac{M-m}{m-1},\quad B=\frac{M-m^2}{m-1},\quad
 r_+=\frac{A+\sqrt{A^2-4B}}2,\quad r_-=B/r_+,
 \quad w=\frac{1-r_-}{r_+-r_-}.
\]
Let $Y$ take values $r_-,r_+$ with probabilities $1-w,w$.
It has moments $\mathbb EY=1$, $\mathbb EY^2=m$, $\mathbb EY^3=M$.
For fixed $e>0$ consider laws $U\ge0$ satisfying
\[
 \mathbb EU=1,\quad\mathbb EU^2\ge m,\quad\mathbb EU^3\le M,
 \qquad \mathbb E\rho(\lambda^2eU)-\lambda pe\ge t
 \quad\hbox{for every }\lambda\ge0.
\]
Then this relaxed class is nonempty if and only if
\begin{equation}\label{eq:flr-domain}
 t\le T_{m,M}(e):=\inf_{\lambda\ge0}
 \{(1-w)\rho(\lambda^2er_-)+w\rho(\lambda^2er_+)-\lambda pe\}.
\end{equation}
When nonempty, its exact remote-payoff supremum is
\begin{equation}\label{eq:flr-value}
 (1-w)f(er_-)+wf(er_+)-pe.
\end{equation}
This is an exact statement about the displayed \emph{relaxed law problem},
not a cubic-realisation theorem. It includes no radial Stein constraints.

\begin{proof}
The functions $f$ and $\rho$ have derivative signs
$f''\le0$, $f'''\ge0$, $f^{(4)}\le0$, and the corresponding signs
for $\rho$. For $\rho$ this follows directly from
\[
 \rho'(u)=\frac1{2\sqrt\pi}\int_0^1e^{-us}
             \sqrt{\frac{1-s}{s}}\,ds;
\]
for $f$ it is the established visibly averaged Bernstein property.
For any function $\phi$ with these three signs, its cubic Hermite
interpolant $H$ at the double nodes $r_-,r_+$ majorises $\phi$ on
the whole half-line. Its cubic coefficient is nonnegative and its
quadratic coefficient nonpositive. Indeed, in Newton form these are
$\phi[r_-,r_-,r_+,r_+]\ge0$ and
$\phi[r_-,r_-,r_+]-(2r_-+r_+)\phi[r_-,r_-,r_+,r_+]\le0$.
Consequently $\mathbb E\phi(U)\le\mathbb E\phi(Y)$.
Apply this simultaneously to $f(eu)$ and every
$\rho(\lambda^2eu)$. If any law satisfies the cuts, $Y$ satisfies
all of them. Conversely $Y$ itself is an admissible relaxed law if
\eqref{eq:flr-domain} holds, and it attains \eqref{eq:flr-value}.
\end{proof}

Thus a new moment-only law LP is unnecessary. The cuts delete part of
the $(e,m,M)$ domain; they do not change the extremiser at a surviving
triple. Adding Stein or tensor constraints can change this conclusion.

\subsubsection{Certified regional remote caps [P,C]}
Assume the project's moment input
\[
 \mathbb E U^2=2+x^2,\qquad
 \mathbb E U^3\le6+9x^2+\frac{4045313}{100000}x^3,
 \qquad 0\le x\le31/4,
 \quad U=|q|^2/e.
\]
Together with \eqref{eq:flr-cut}, it gives
\[
\begin{array}{c|c|c}
 \text{condition}&\mathbb E f(|q|^2)-pe\text{ is less than}
      &\text{improvement over }1.3739635\\\hline
 t\ge0.847&1.3739601&0.0000034\\
 t\ge0.849&1.3739481&0.0000154\\
 t\ge0.851&1.3739276&0.0000359
\end{array}
\]
All displayed decimals are exact rational numbers. These estimates allow
all shared sectors in the original response, but estimate only its
remote component. They require no independent Gaussian buffer.

\paragraph{Certificate mechanism.}
On a rational kurtosis band $[x_-,x_+]$, use
$m=2+x_-^2$ and $M=6+9x_+^2+(4045313/100000)x_+^3$.
Fix $\lambda=39/25$. The function
\[
 H(e)=(1-w)\rho(\lambda^2er_-)+w\rho(\lambda^2er_+)-\lambda pe
\]
is strictly decreasing for all $e\ge0$, because
$H'(e)\le\lambda^2\sqrt\pi/4-\lambda p<0$.
An enclosed value $H(e_*)<t$ therefore certifies $e<e_*$,
with no unchecked energy tail. The remote objective in
\eqref{eq:flr-value} is concave in $e$. A positive forward secant at
$e_*$ proves that its maximum on $[0,e_*]$ is its endpoint value.
Otherwise the existing four-point concavity enclosure bounds its
maximum on the entire energy half-line.

The fresh replay checks 43, 39 and 31 kurtosis bands respectively.
It recomputes every radial integral using Arb balls, Jensen/chord
bounds on each radial cell, and an explicit Gaussian tail.
Stored evaluation values are not used for acceptance. Thus [C] here
means directed interval computation, not floating optimisation or a
finite radial cutoff. The inputs are the existing derivative-sign
theorem and project moment bound, plus the proved response cut.

\subsubsection{A certified limit on these scalar cuts [P,C]}
For $m=2$, $M=6$ the extremal law has nodes $2\pm\sqrt2$ and
upper probability $(2-\sqrt2)/4$. Since $E\sim\mathrm{Exp}(1)$
has the same first three moments, the two-node theorem gives, for
every $\lambda\ge0$,
\[
 \mathbb E\rho(\lambda^2eY)\ge
 \mathbb E\rho(\lambda^2eE)
 =\frac{\sqrt\pi}{2}\sqrt{1+\lambda^2e}.
\]
Consequently this law passes \emph{every} scalar Rice cut whenever
\begin{equation}\label{eq:flr-survivor}
 t^2+p^2e\le\pi/4.
\end{equation}
This follows either by minimising the last display minus $\lambda pe$,
or by the Euclidean Cauchy--Schwarz inequality. At
\[
 t=1057/1250,\qquad e=161/4000,
\]
\eqref{eq:flr-survivor} is strict and the fresh directed computation gives
\[
 \mathbb E f(eY)-pe>1.37396307.
\]
It also certifies $S<1.37395354$ on the entire Gaussian-variance
half-line. Hence this surviving abstract law has payoff
$>S+9.53\cdot10^{-6}$. At $t=0.84$, choosing $e=651/16000$
gives payoff $>1.37396328$ while still satisfying all cuts.
The law also satisfies the existing radius bound:
$(\mathbb E\sqrt Y)^2=(2+\sqrt2)/4<0.9718^2$.

Therefore, relative to the old remote cap $1.3739635$, these moment
constraints and \emph{all} scalar Rice cuts can improve the peak cap
by less than $4.3\cdot10^{-7}$, and the left-end cap by less than
$2.2\cdot10^{-7}$. This is not a counterexample to K and does not
claim feasibility for Stein cuts or actual cubic responses. Strict
positive excess kurtosis does not remove the obstruction: continuity
allows sufficiently small $x>0$, because the support inequality above
has a strict margin at the stated energy.

\subsubsection{A joint response-projection constraint [P]}
There is a stronger consequence of feasibility which keeps the tensor
information discarded by a scalar law. Put $\bar k=\mathbb E[k\mid X]$
and
\[
 r_\lambda(q)=2\lambda q\rho'(\lambda^2|q|^2),\qquad
 D_\lambda=\mathbb E\rho(\lambda^2|q|^2)-t-\lambda pe.
\]
Then, for every $\lambda\ge0$,
\begin{equation}\label{eq:flr-response-stability}
 \boxed{\quad
 \|\bar k-r_\lambda(q)\|_2^2\le\sqrt\pi\,D_\lambda.
 \quad}
\end{equation}
No moment bound or independence among remote fields is needed.

\begin{proof}
Every remote $v\in L^2(X;\mathbb C)$ obeys
\[
 t+\operatorname{Re}\langle v,\bar k\rangle
 \le\mathcal F(v):=\mathbb E\rho(|v|^2).
\]
This is the same two-response support argument, without requiring
$v$ to belong to third chaos. Treat $L^2(X;\mathbb C)$ as a real
Hilbert space. The gradient is
$\nabla\mathcal F(v)=2v\rho'(|v|^2)$, and
\[
 D^2\mathcal F(v)[u,u]
 =\mathbb E\{2\rho'(|v|^2)|u|^2
       +4\rho''(|v|^2)(\operatorname{Re}\bar vu)^2\}
 \le\frac{\sqrt\pi}{2}\|u\|_2^2.
\]
Convexity follows directly from the Gaussian average of the norm.
The upper Hessian bound gives a global quadratic majorant, also for
unbounded $L^2$ arguments by approximation. At $v=\lambda q$,
compare that majorant at $v+u$ with the support lower bound there:
\[
 D_\lambda\ge
 \operatorname{Re}\langle u,\bar k-r_\lambda(q)\rangle
       -\frac{\sqrt\pi}{4}\|u\|_2^2.
\]
Choosing $u=2(\bar k-r_\lambda(q))/\sqrt\pi$ proves
\eqref{eq:flr-response-stability}. The identity
$\langle q,\bar k\rangle=pe$ is the only use of $F=\varepsilon Qk$
in evaluating its support gap.
\end{proof}

\paragraph{Both conditional responses can be retained.}
In fact, with $\bar h=\mathbb E[h\mid X]$, the stronger statement is
\begin{equation}\label{eq:flr-both-responses}
 \boxed{\quad\|\bar h\|_2^2+
 \|\bar k-r_\lambda(q)\|_2^2\le\sqrt\pi D_\lambda.\quad}
\end{equation}
To see this, use the two-variable support functional
\[
 \mathcal F(u,v)=\tfrac12\mathbb E\{\rho(|u+v|^2)+\rho(|u-v|^2)\}
 \ge t+\operatorname{Re}\langle u,\bar h\rangle
          +\operatorname{Re}\langle v,\bar k\rangle.
\]
Its real Hessian on the product Hilbert space is again bounded by
$\sqrt\pi/2$, by the parallelogram identity. Its gradient at
$(0,\lambda q)$ is $(0,r_\lambda(q))$. The same square completion
proves \eqref{eq:flr-both-responses}.

Let $P_X,Q_X$ denote the remote first and $(2,1)$ chaos projections.
Since $Q_X\bar k=pq$ and the two ranges are orthogonal,
\begin{equation}\label{eq:flr-joint}
 \boxed{\quad
 \|pq-Q_Xr_\lambda(q)\|_2^2
 +\|P_X\bar k-P_Xr_\lambda(q)\|_2^2
 \le\sqrt\pi D_\lambda.\quad}
\end{equation}
The additional term $\|Q_X\bar h\|_2^2$ may be retained on its left.
In particular $\|pq-Q_Xr_\lambda(q)\|_2^2>\sqrt\pi D_\lambda$
is a certificate of infeasibility, even if every scalar cut passes.
If
\[
 a=\|P_Xr_\lambda(q)\|_2,\qquad
 b^2=\|pq-Q_Xr_\lambda(q)\|_2^2,\qquad
 \tau=\sqrt\pi D_\lambda-b^2\ge0,
\]
the original row's retained negative square satisfies
\begin{equation}\label{eq:flr-penalty}
 -\|Pk\|_2^2\le-(a-\sqrt\tau)_+^2.
\end{equation}
This is a lower estimate of an existing penalty, not an additional
penalty to subtract twice from an already optimised response dual.

\paragraph{Finite implementation.}
One may replace $P_X$ and $Q_X$ in \eqref{eq:flr-joint} by arbitrary
finite orthonormal test subspaces of those respective chaoses.
All coefficients are then finite integrals of the same explicit
profile $r_\lambda(q)$. The resulting norm sums give a safe weaker
test and penalty. In contrast to scalar moment cuts, these retain
the cubic tensor through its nonlinear first- and third-chaos
projections. Quantifying their gain in a full all-sector row remains
open; this package does not claim that gain.

\paragraph{Next necessary input.} Section~\ref{sec:coupled-projection-certificate} now carries out that test on a genuine nine-coordinate cubic with normalized
fourth moment $593/100$. Every scalar Rice cut passes, but a noncollinear cubic support test excludes it. The same test gives a necessary bound on two projections of
every remote cubic, leaving the orthogonal component unrestricted. A nearby two-coefficient region forces $\|Pk\|^2>10^{-5}$. The remaining task is to quantify these
gains on the high-payoff tensors in a complete original-response row, not to apply the cuts to an auxiliary price-layer law.

\subsection{A coupled cubic that passes every scalar support cut but fails response feasibility}
\paragraph{Provenance (IDE note, 21 October).} Lower-bound GPT coupled-projection package (\code{conjK-coupled-projection-certificate-20261021.zip}, archived in
\pathref{gpt-reviews/coupled-projection-20261021/}). IDE replays \code{python -I verify\_all.py} and \code{--fine} (pinned \code{python-flint 0.9.0}, unrelated directory):
ALL CLAIMS PASS in 3.1\,s and 5.6\,s, with fresh directed integrals and tails. IDE check: $\E|J|^2=\frac14\E|X|^2\sum_j\E|Y_j|^2\E|W_j|^2=1$ and $H\perp J$ (since
$\E[\bar Y_jW_j]=0$), so $\|H+3J\|^2=10$. This is a constraint on original responses, not a raw-K family; the conditional $\|Pk\|^2$ bound bounds an existing negative
term and must not be subtracted twice.

\label{sec:coupled-projection-certificate}
\paragraph{Status.}
The analytic statements below are proved [P]. The displayed rational
margins are certified by fresh directed integration [C]. This is a
theorem about original bounded responses, not a proof of raw K or a new
Grothendieck endpoint. No new moment or chaos conjecture is assumed.

\subsubsection{Setting and the finite-dimensional cubic}
Let $Z$ be the visible proper complex Gaussian, independent of the
remote standard proper Gaussians $X,Y_1,\ldots,Y_4,W_1,\ldots,W_4$.
Write
\[
 H=\frac{X(|X|^2-2)}{\sqrt2},\qquad
 V=\frac12\sum_{j=1}^4Y_j\overline W_j,\qquad J=XV.
\]
The remote cubics $H,J$ are orthonormal in $H_{2,1}$, and $X$ is
orthogonal to both. Put
\[
 p=2500/1891,\qquad t_0=1057/1250,\qquad
 q_e=\sqrt{e/10}\,(H+3J),\qquad \lambda_0=42/25.
\]
In particular $\|q_e\|_2^2=e$.
If $T=|X|^2$ and $Y=\sum|Y_j|^2$, then $T\sim\mathrm{Exp}(1)$,
$Y\sim\mathrm{Gamma}(4,1)$ independently, and conditionally on
$X,Y_1,\ldots,Y_4$, the $dJ$ term is proper Gaussian of variance
$d^2TY/4$. Therefore
\[
 \mathbb E|cH+dJ|^4=44c^4+16c^2d^2+5d^4,
 \qquad \frac{\mathbb E|q_e|^4}{e^2}=\frac{593}{100}.
\]
This is a genuine finite cubic in nine coordinates, with a Wick
self-interaction and a Gaussian-factor term sharing $X$. Its kurtosis
lies inside the application's range $[2.06,6]$.

Here $P=P_{1,0}$ and $Q=P_{2,1}$ denote the full orthogonal chaos
projections, and $\varepsilon=1891/2500$. For the original responses let $|f_\pm|\le1$,
$h=(f_++f_-)/2$, $k=(f_+-f_-)/2$, $Ph=tZ$, and
$q=\mathbb E[\varepsilon Qk\mid\text{remote}]$.
Every visible sector is unrestricted. Set
\[
 \rho(u)=\mathbb E|Z+\sqrt u|,\qquad R=\sqrt\pi/2,
 \quad \mathcal F(v)=\mathbb E\rho(|v|^2),\quad
 r(v)=2v\rho'(|v|^2).
\]
Every remote test $v\in H_{2,1}$ satisfies
\begin{equation}\label{eq:jpc-support}
 t+p\operatorname{Re}\langle v,q\rangle\le\mathcal F(v).
\end{equation}
Scalar Rice cuts restrict $v$ to $\lambda q$. The certificate below
uses a test with a different ratio of the $H$ and $J$ coefficients.
To see the support bound directly, for any remote $u\in L^2$ use
\[
 \operatorname{Re}(\overline Z h+\overline u k)
 \le\tfrac12\{|Z+u|+|Z-u|\}.
\]
Its expectation is $t+\operatorname{Re}\langle u,\bar k\rangle$
on the left and $\mathcal F(u)$ on the right, where
$\bar k=\mathbb E[k\mid\text{remote}]$. Gaussian conditional expectation
gives $Q_X\bar k=pq$. This proves \eqref{eq:jpc-support} without
any restriction on the remaining projections of the responses.

\subsubsection{Main certified conclusions [P,C]}
Let $e_0=31/800$ and $c_0=\sqrt{31/8000}$.

\paragraph{Every scalar cut passes at the excluded cubic.}
At $q=q_{e_0}$ and $t=t_0$,
\begin{equation}\label{eq:jpc-all-scalar}
 \boxed{\quad \mathcal F(\lambda q)-t_0-\lambda pe_0
     >\frac1{100000}\quad\hbox{for every }\lambda\ge0.\quad}
\end{equation}
This is a continuum statement, including both unbounded ends of the
$\lambda$ domain. It is not inferred from a grid of multipliers.

\paragraph{Nevertheless the response configuration is impossible.}
Define the positive real algebraic coefficients
\[
 a_0=\lambda_0c_0+\frac7{1000},\qquad
 b_0=3\lambda_0c_0-\frac{23}{10000},\qquad v_*=a_0H+b_0J.
\]
The replay proves
\begin{equation}\label{eq:jpc-separated}
 \mathcal F(v_*)-t_0-p\langle v_*,q_{e_0}\rangle
 <-\frac1{100000},\qquad \|v_*\|_2<\frac13.
\end{equation}
Hence there are no original responses at $t\ge t_0$ with remote
projection $q_{e_0}$, despite \eqref{eq:jpc-all-scalar}.

\paragraph{A universal constraint on two cubic coefficients.}
For \emph{any} actual remote cubic $q$, with arbitrary additional remote
coordinates and dependencies, put $c=\langle H,q\rangle$ and
$d=\langle J,q\rangle$ (the inner product is linear in its second slot).
Then, at $t\ge t_0$,
\begin{equation}\label{eq:jpc-universal-cut}
 \boxed{\quad a_0|c|+b_0|d|
 <(a_0+3b_0)c_0-\frac1{100000p}.\quad}
\end{equation}
There is \emph{no size restriction on the component of $q$ orthogonal
to $H,J$}. In particular no original response can have coefficients
satisfying
\begin{equation}\label{eq:jpc-halfspace}
 a_0|c|+b_0|d|\ge(a_0+3b_0)c_0.
\end{equation}
These statements hold for any nine orthonormal remote Gaussian directions
used to form $H,J$, not only for fixed coordinates.
There is also a direct row-parameter version. Define the fixed algebraic
constant $C_*=t_0+p(a_0+3b_0)c_0-10^{-5}$. For any
$\varepsilon'>0$, responses with $Ph=tZ$ and actual remote projection
$q'=\mathbb E[\varepsilon'Qk\mid\text{remote}]$ satisfy
\begin{equation}\label{eq:jpc-any-row}
 t+\frac{a_0|\langle H,q'\rangle|+b_0|\langle J,q'\rangle|}
              {\varepsilon'}<C_*.
\end{equation}
The same certified test therefore applies across row parameters. It does
not involve the auxiliary visible profile $L$ at all.

\begin{proof}
Choose the phases of the coefficients of $v_*$ to align its pairing with
$c,d$. Rotating $X$ and the common phase of the $Y_j$ variables permits
the phases of $H,J$ to be set independently, so the distribution of the
resulting test, and hence $\mathcal F(v_*)$, is unchanged. The support
pairing is exactly $a_0|c|+b_0|d|$. Its pairing with every orthogonal
component is zero. Apply \eqref{eq:jpc-support} and
\eqref{eq:jpc-separated}. Gaussian unitary invariance allows arbitrary
orthonormal remote directions.
\end{proof}

As a convenient neighbourhood consequence, if complex $c,d$ obey
\eqref{eq:jpc-halfspace}, there are no original responses with
\begin{equation}\label{eq:jpc-robust}
 q=cH+dJ+q_\perp+\Delta,\qquad
 q_\perp\perp H,J,\qquad \|\Delta\|_2\le1/50000.
\end{equation}
The cubic $q_\perp$ is unrestricted in size and dependence. The remaining
support violation exceeds $10^{-6}$, since its decrease is at most
$p\|v_*\|\|\Delta\|<p/150000$ and
$1/100000-p/150000>1/1000000$.
In particular this excludes every $sq_{e_0}+q_\perp+\Delta$, $s\ge1$,
under these same conditions.

\paragraph{A forced penalty at a nearby cubic.}
Let $e_1=38739/10^6$. If an original response at $t\ge t_0$ has remote
projection $q_{e_1}$, then
\begin{equation}\label{eq:jpc-penalty-center}
 \|Pk\|_2^2>1/75000.
\end{equation}
More robustly, if its remote projection is $q_{e_1}+\Delta$, where
$\Delta$ is any remote cubic with $\|\Delta\|_2\le10^{-6}$, then
\begin{equation}\label{eq:jpc-penalty-ball}
 \boxed{\quad \|Pk\|_2^2>10^{-5}.\quad}
\end{equation}
The central cubic passes every scalar cut with gap $>2\cdot10^{-5}$.
More generally, write $c=\langle H,q\rangle$,
$d=\langle J,q\rangle$, and $c_1=\sqrt{e_1/10}$. The same bound
\eqref{eq:jpc-penalty-ball} holds whenever
\begin{equation}\label{eq:jpc-penalty-projections}
 (|c|-c_1)^2+(|d|-3c_1)^2\le10^{-12},
\end{equation}
with \emph{every orthogonal remote cubic component unrestricted}.
If $|c|=c_1$ and $|d|=3c_1$, the stronger bound $1/75000$ holds.
Thus the result concerns a neighbourhood of two coefficient magnitudes,
not a small ball around a whole high-dimensional tensor.

These penalty assertions are conditional on a response existing;
no such response is constructed here. The quantity $-\|Pk\|_2^2$
is already in the exact row identity, and cannot be deducted twice.

\subsubsection{Why the certification needs only two radial variables [P]}
For $q=cH+dJ$, $c,d\ge0$, and any $\lambda\ge0$, put
\[
 A=Y/4,\quad Z_0(T)=T(T-2)^2/2,\quad
 s=1+\lambda^2d^2TA,\quad u=\lambda^2c^2Z_0(T)/s,
 \quad K_j(u)=e^{-u/2}I_j(u/2).
\]
The following expectations are over the independent $T,Y$ only:
\begin{align}
 \mathcal F(\lambda q)
 &=R\mathbb E\big[\sqrt s\{(1+u)K_0(u)+uK_1(u)\}\big],\label{eq:jpc-A}\\
 r_H:=\langle H,r(\lambda q)\rangle
 &=\lambda cR\mathbb E\left[\frac{Z_0(T)}{\sqrt s}(K_0+K_1)\right],\label{eq:jpc-H}\\
 r_J:=\langle J,r(\lambda q)\rangle
 &=\lambda dR\mathbb E\left[\frac{TA}{\sqrt s}K_0\right],\label{eq:jpc-J}\\
 r_X:=\langle X,r(\lambda q)\rangle
 &=\lambda cR\mathbb E\left[\frac{T(T-2)}{\sqrt{2s}}(K_0+K_1)\right].\label{eq:jpc-X}
\end{align}
All these coefficients are real. The first formula averages the Gaussian
$W_j$ coordinates together with the visible unit Gaussian in $\rho$.
Differentiating it with respect to $c,d$ yields
$\partial_c\mathcal F(\lambda q)=\lambda r_H$ and
$\partial_d\mathcal F(\lambda q)=\lambda r_J$; here
$\rho(u)-2u\rho'(u)=RK_0(u)$.
Conditioning the phase response directly gives \eqref{eq:jpc-X}.
In particular, if
$D(\lambda)=\mathcal F(\lambda q)-t_0-\lambda pe$, then
\begin{equation}\label{eq:jpc-derivative}
 D'(\lambda)=cr_H+dr_J-pe.
\end{equation}
No angular quadrature, large-dimension truncation, or Gaussian-limit
replacement is used.

\subsubsection{A global test of all scalar multipliers [P]}
The unit Gaussian norm has the Hessian bounds
\begin{equation}\label{eq:jpc-hessians}
 R e^{-|w|^2}I\ \preceq\ D^2_w\mathbb E|Z+w|\ \preceq\ RI.
\end{equation}
For example, use
$|w|=\frac14\int_0^{2\pi}|\operatorname{Re}(e^{-i\theta}w)|\,d\theta$.
The real Gaussian projection has density $e^{-x^2}/\sqrt\pi$, so its
smoothed absolute value has second derivative $2e^{-x^2}/\sqrt\pi$.
Integrating a directional square gives both inequalities in
\eqref{eq:jpc-hessians}.

If $M_4=\mathbb E|q|^4$ and $e=\|q\|_2^2>0$, Jensen under the
$|q|^2$-size-biased law gives
\[
 D''(\lambda)\ge R\mathbb E[|q|^2e^{-\lambda^2|q|^2}]
       \ge Re\,e^{-\lambda^2M_4/e}.
\]
For both certified cubics, the replay checks
\[
 Re\exp\{-9(593/100)e\}>1/250=:\kappa.
\]
Thus $D$ is $\kappa$-strongly convex on $[0,3]$.
It also checks $D'(\lambda_0)+\kappa(3-\lambda_0)>0$;
global convexity then makes $D$ increasing beyond $3$. Consequently
\begin{equation}\label{eq:jpc-all-lambda-check}
 \inf_{\lambda\ge0}D(\lambda)
 \ge D(\lambda_0)-\frac{D'(\lambda_0)^2}{2\kappa}.
\end{equation}
Directed evaluation gives a lower bound greater than $1.1943\cdot10^{-5}$
at $e_0$ and $2.4167\cdot10^{-5}$ at $e_1$.

\subsubsection{The noncollinear test and the penalty [P,C]}
From the upper Hessian bound in \eqref{eq:jpc-hessians}, for any remote
$v,w\in L^2$,
\[
 \mathcal F(v+w)\le\mathcal F(v)
       +\operatorname{Re}\langle w,r(v)\rangle+\frac R2\|w\|^2.
\]
At $v=\lambda_0q_{e_0}$ choose
$w=(7/1000)H-(23/10000)J$. This proves
\eqref{eq:jpc-separated} by a single directed comparison of
\[
 D(\lambda_0)+(r_H-pc_0)\frac7{1000}
 -(r_J-3pc_0)\frac{23}{10000}
 +\frac R2\left\{\left(\frac7{1000}\right)^2+
                  \left(\frac{23}{10000}\right)^2\right\}
 <-10^{-5}.
\]

For completeness, the response projection constraint used for the penalty
also follows from this same upper Hessian bound. Put $\bar k=\mathbb E[k\mid X,Y,W]$.
The support inequality for every remote $u\in L^2$ reads
$t+\operatorname{Re}\langle u,\bar k\rangle\le\mathcal F(u)$.
The quadratic majorant at a fixed $v$, optimised in its increment, gives
\[
 \|\bar k-r(v)\|^2\le2R\{\mathcal F(v)-t-
                           \operatorname{Re}\langle v,\bar k\rangle\}.
\]
Take $v=\lambda_0q_{e_1}$ and let the actual remote projection be
$q_{e_1}+\Delta$, $\|\Delta\|\le\delta$.
In the orthonormal directions $H,J,X$, write
\[
 a=|r_X|,\qquad b=\sqrt{(pc-r_H)^2+(pd-r_J)^2},\qquad
 D=\mathcal F(v)-t_0-\lambda_0pe_1.
\]
Bessel's inequality and the perturbation norm give
\[
 \|Pk\|^2\ge
 \left(a-\sqrt{\left[\sqrt\pi\{D+p\|v\|\delta\}
                    -(b-p\delta)_+^2\right]_+}\right)_+^2,
\]
whenever the response is feasible; a negative bracket already excludes it.
The replay checks the strict lower bounds $1/75000$ at $\delta=0$
and $1/100000$ at $\delta=10^{-6}$. Only the projections of
$\Delta$ onto $H,J$ entered either the support pairing or Bessel's
inequality. Thus one may replace $\|\Delta\|\le\delta$ by
$|\langle H,\Delta\rangle|^2+|\langle J,\Delta\rangle|^2\le\delta^2$.
Aligning coefficient phases as above proves
\eqref{eq:jpc-penalty-projections}. Additional remote coordinates,
unbounded orthogonal components, and arbitrary dependence do not affect
the argument.

\subsubsection{Directed integration and tails [C]}
The replay integrates \eqref{eq:jpc-A}--\eqref{eq:jpc-X} on
$0\le T\le36$, $0\le Y\le48$ against
$e^{-T-Y}Y^3/6\,dT\,dY$.
Arb's complex analytic integrator supplies directed enclosures, not
estimated floating quadrature errors. On every complex argument ball
used for an analytic bound, the implementation verifies $\operatorname{Re}s>0$.
This fixes a holomorphic square-root branch and excludes poles.

For wide complex balls it avoids the interval cancellation between
$e^{-u/2}$ and $I_j(u/2)$ by using the positive Laplace representations
\[
 K_0(u)={}_1F_1(1/2;1;-u),\qquad
 K_0(u)+K_1(u)={}_1F_1(1/2;2;-u).
\]
Both are transforms of probability measures on $[0,1]$. If
$E_u=\exp(\max(0,-\operatorname{Re}u))$, their moduli are at most
$E_u$, while
$|\rho(u)/R|\le1+|u|E_u/2$.
These bounds justify the complex range enclosures used by the integrator.
At narrow balls the entire Bessel formulas are evaluated directly.

On the omitted real tails, nonnegative polynomial majorants for the
four integrands before multiplying by the probability density are
\[
\begin{array}{c|c}
 \mathcal F&3+4T+T^2+Y/4\\
 |r_X|&T^2+2T\\
 r_H&T^3/2+2T^2+2T\\
 r_J&TY/4.
\end{array}
\]
They follow from $K_0,K_0+K_1\le1$, $s\ge1$, and the triangle
inequality for the norm. For example, at the certified parameters
$\lambda c<1$ and $\lambda dR<1$, so the support integrand is at most
\[
 R\sqrt s+\lambda c\sqrt{Z_0(T)}
 \le 1+\tfrac12(T+Y/4)+\tfrac12(T^2+3T+2),
\]
which is smaller than its displayed majorant.
Exact integer-shape incomplete gamma formulas
bound the union of $T>36$ and $Y>48$. All four tail errors are below
$10^{-10}$; the largest is below $6.541\cdot10^{-12}$.
No stored function value or floating optimisation output is trusted.

\paragraph{What this changes.}
This is the first test in this continuation where a \emph{genuine}
cubic in the needed kurtosis range passes all scalar Rice tests and
is excluded by a tensor-dependent response test. The same test constrains
two projections of every cubic, with the entire orthogonal component
unrestricted. Nearby in these two coefficients it gives a conditional
$10^{-5}$ deduction from an existing row penalty.
It does not yet prove that high-payoff candidates across the complete
row domain receive such a deduction. A full all-sector optimisation,
then the profile envelope, is still necessary for a new endpoint.

\subsection{All-sector response feasibility: an effective first moment and a joint profile constraint}
\label{sec:sector-rice-feasibility}
\paragraph{Provenance (IDE note, 21 October).} Lower-bound GPT sector-feasibility package (\code{conjK-all-sector-response-constraints-20261021.zip}, archived in
\pathref{gpt-reviews/sector-feasibility-20261021/}). IDE replay \code{python3 -I verify\_all.py}: ALL CLAIMS PASS (exact rational algebra, 0.05\,s). IDE re-derivation \st{P}
of \eqref{eq:sector-ellipsoid}: $|Z+U|+|Z-U|\le2|Z|+|U|^2/|Z|$; the weights $\E[w_k(T)/\sqrt T]/R$ are $2,\frac32,\frac34,1,1,\frac78$ (via $\Gamma(\frac12),\Gamma(\frac32),\Gamma(\frac52),
\Gamma(\frac72)$); cross terms vanish by the visible phase, by remote orthogonality, or by $\E Y=0$; and optimising each test scale $\lambda_k=p/(c_kR)$ gives the
coefficients $1/(2c_k)$.

\paragraph{Status and scope.}
The dimension-free statements below are analytic proofs [P]. Their few
rational polynomial bounds have a standard-library exact replay [C].
The candidate diagnostics are floating calculations [F]. No raw-K theorem
or new Grothendieck endpoint is claimed.

Let $P=P_{1,0}$, $Q=P_{2,1}$, $\varepsilon>0$, and $p=1/\varepsilon$.
For bounded responses $|f_\pm|\le1$, put
\[
 h=(f_++f_-)/2,\qquad k=(f_+-f_-)/2,\qquad Ph=tZ,
 \qquad F=\varepsilon Qk.
\]
The visible standard proper Gaussian $Z$ is independent of all remote
coordinates $X$. Write the full cubic as
\[
 F=q+(|Z|^2-1)r_A+\frac{Z^2}{\sqrt2}r_B
       +ZY+\overline ZW+\nu H(Z),\qquad
 H(Z)=Z(|Z|^2-2)/\sqrt2.
\]
Here $q,r_A,r_B,Y,W$ are arbitrary jointly defined remote chaoses of
bidegrees $(2,1),(1,0),(0,1),(1,1),(2,0)$, respectively.
Their dependencies and phases are unrestricted. The scalar $\nu$ is
arbitrary. Define
\[
 r_\parallel=\frac{\sqrt2r_A+\overline{r_B}}{\sqrt3},\qquad
 r_\perp=\frac{r_A-\sqrt2\overline{r_B}}{\sqrt3},\qquad
 E_{\rm ev}=\|Y\|_2^2+\|W\|_2^2+|\nu|^2.
\]
This is a real orthogonal change of variables: the sum of the two squared
norms is $\|r_A\|^2+\|r_B\|^2$.

Put $R=\sqrt\pi/2$, $\rho(u)=\mathbb E|Z+\sqrt u|$, and, for any
\emph{remote} cubic test $v$, set
\[
 \mathcal D_t(v;q)=\mathbb E\rho(|v|^2)-t-p\operatorname{Re}\langle v,q\rangle.
\]
The support inequality makes this nonnegative for actual responses.

\subsubsection{An exact energy ellipsoid halves the full cubic domain [P,C]}
There is a particularly simple all-sector consequence of retaining the
first-moment constraint:
\begin{equation}\label{eq:sector-ellipsoid}
 \boxed{\quad
 \frac{\|q\|^2}{4}+\frac{\|r_A\|^2}{3}
 +\frac{2\|r_B\|^2}{3}
 +\frac{\|Y\|^2+\|W\|^2}{2}+\frac{4|\nu|^2}{7}
 \le R\varepsilon^2(R-t).\quad}
\end{equation}
In particular $t\le R$ for any response. Together with the scalar
Rice cut at zero's tangent, this gives the closed-form bound
\begin{equation}\label{eq:sector-closed-total}
 \boxed{\quad \|F\|^2\le\frac72 R\varepsilon^2(R-t).\quad}
\end{equation}
Using the slightly sharper rational Rice tangent below gives
\begin{equation}\label{eq:sector-refined-total}
 \|F\|^2\le\varepsilon^2
   \{3R(R-t)+B_0(A_0-t)\}.
\end{equation}
For $\varepsilon=1891/2500$, the exact checker proves
\[
\begin{array}{c|c|c}
 t\text{ lower bound}&\text{new strict cap on }\|F\|^2
 &\text{old }\varepsilon^2(1-t^2)\text{ cap, approximately}\\\hline
 0.840&0.081884&0.168438299\\
 0.8456&0.071981&0.163037655\\
 0.851&0.062431&0.157795905
\end{array}
\]
These are bounds on the \emph{whole} actual cubic, including every shared
sector, not just on the remote component. They do not imply a payoff cap.

\begin{proof}
Take an arbitrary full cubic test $U$. The support inequality and the
global norm majorant at $Z$ give
\[
 t+p\operatorname{Re}\langle U,F\rangle
 \le\tfrac12\mathbb E(|Z+U|+|Z-U|)
 \le R+\tfrac12\mathbb E\frac{|U|^2}{|Z|}.
\]
The weighted quadratic form is diagonal on the six cubic sectors.
Angular orthogonality removes most cross terms. The only remaining
potential cross terms are remote versus $A$, and $Y$ versus diagonal;
they vanish because their remote coefficient fields lie in orthogonal
chaoses (and $\mathbb EY=0$). Divided by $R$, its six coefficients are
\[
 2,\quad 3/2,\quad3/4,\quad1,\quad1,\quad7/8.
\]
For example the $A$ coefficient is
$\mathbb E[(T-1)^2/\sqrt T]/R=3/2$, and the diagonal coefficient is
$\mathbb E[T(T-2)^2/(2\sqrt T)]/R=7/8$.
Completing six independent squares proves
\eqref{eq:sector-ellipsoid}. This argument requires no independence
between the remote coefficient fields.

Let $e=\|q\|^2$ and $C=R\varepsilon^2(R-t)$. Every visible coefficient
in \eqref{eq:sector-ellipsoid} is at least $1/3$, hence
$\|F\|^2\le3C+e/4$.
The global tangent $\rho(u)\le R+(R/2)u$ gives $e\le2C$ by
optimizing $\lambda$ in the scalar cut. This proves
\eqref{eq:sector-closed-total}. The refined remote bound
$e\le4\varepsilon^2B_0(A_0-t)$ proves
\eqref{eq:sector-refined-total}.
\end{proof}

\subsubsection{Two universal inequalities [P]}
For every such test $v$,
\begin{equation}\label{eq:sector-effective-moment}
 \boxed{\quad
 \frac{3\|r_\perp\|^2+2E_{\rm ev}}{8\varepsilon^2}
 \le\mathcal D_t(v;q).\quad}
\end{equation}
It holds with \emph{all sectors present simultaneously}, in every
finite dimension, at every amplitude, without a fresh buffer and without
stationarity. In particular it controls the interaction of even sectors
with arbitrary shared $A/B$ fields, as a response-feasibility constraint.
It is not a comparison of raw payoffs after adding those fields.

There is a complementary inequality retaining the parallel content.
Define the bounded positive kernel
\[
 G(u)=e^{-u/2}\{I_0(u/2)-I_1(u/2)\}
 =\frac2\pi\int_0^1e^{-us}\sqrt{\frac{s}{1-s}}\,ds,
 \qquad 0<G(u)\le1,
\]
and the remote first-chaos projection
\[
 a(v)=-\frac R2 P_X\{vG(|v|^2)\}.
\]
Then
\begin{equation}\label{eq:sector-joint-profile}
 \boxed{\quad
 \left\|p r_\parallel-\sqrt{\frac32}\,a(v)\right\|^2
       +p^2\|r_\perp\|^2+p^2E_{\rm ev}
 \le4\mathcal D_t(v;q).\quad}
\end{equation}
The profile is averaged over the visible Gaussian exactly before any
Cauchy--Schwarz bound. Since $P_Xv=0$, it may equivalently be evaluated as
$a(v)=(R/2)P_X\{v(1-G(|v|^2))\}$, which removes the linear term.
No estimate of a third or fifth derivative of the old $L$-profile occurs.

The two inequalities are valid simultaneously. Their left sides must
\emph{not} be added while retaining the same right side: they are two
optimizations of a shared support budget. The same caution applies to
combining them with a separate remote-projection certificate. Section~\ref{sec:joint-response-budget} proves which terms \emph{can} be combined, using uniformity of the
transverse/even support charge over every remote shift; it also gives a separate optimized parallel-profile envelope.

\subsubsection{Consequences for the existing remote cuts [P,C]}
Define
\[
 t_{\rm eff}=t+\frac{3\|r_\perp\|^2+2E_{\rm ev}}{8\varepsilon^2}.
\]
Every remote support cut can be strengthened simply by replacing $t$ by
$t_{\rm eff}$:
\begin{equation}\label{eq:sector-all-tests}
 t_{\rm eff}+p\operatorname{Re}\langle v,q\rangle
       \le\mathbb E\rho(|v|^2).
\end{equation}
This includes noncollinear tensor tests. At $v=\lambda q$, it gives a
linear constraint on the law of $|q|^2$ and the two sector energies:
\[
 \mathbb E[\rho(\lambda^2|q|^2)-\lambda p|q|^2]
 \ge t_{\rm eff}\qquad(\lambda\ge0).
\]
This substitution is a necessary constraint, not an assertion that a new
response with first projection $t_{\rm eff}Z$ has been constructed.

The existing certified global tangent
\[
 \rho(u)\le A_0+B_0u,\qquad
 A_0=886762780/10^9,\quad B_0=432306903/10^9
\]
therefore gives the global coupled energy bound
\begin{equation}\label{eq:sector-energy-cut}
 \boxed{\quad
 \|q\|^2+\frac{3B_0}{2}\|r_\perp\|^2+B_0E_{\rm ev}
 \le4\varepsilon^2B_0(A_0-t).\quad}
\end{equation}
The checker recomputes the tangent using rational alternating series.
At $\varepsilon=1891/2500$, the right side is strictly below
$0.046266$, $0.040726$, and $0.035383$ for
$t\ge0.84$, $0.8456$, and $0.851$, respectively.
For scale, $\|r_\perp\|=0.012$ consumes more than $94\cdot10^{-6}$
of the first-moment support budget. This is \emph{not} a deduction of
that size from the norm cap.

The old rotation identity pays for $r_\parallel$ because
$2\beta r_A+\sqrt2\beta\overline{r_B}=\sqrt6\beta r_\parallel$.
Equation~\eqref{eq:sector-effective-moment} controls its orthogonal
combination through feasibility, including even and diagonal interactions.
Equation~\eqref{eq:sector-joint-profile} further ties the parallel field
to one explicit remote profile. These are quantitative structural
constraints; their combination with the full row objective is still an
open optimization problem.

\subsubsection{The visible Gaussian calculation [P]}
For a fixed remote shift $w$, define
\[
 k_w(Z)=\tfrac12\{\operatorname{ph}(Z+w)-\operatorname{ph}(Z-w)\}.
\]
It is even under $Z\mapsto-Z$. Let
$r(w)=\mathbb E_Z\operatorname{ph}(Z+w)=wf(|w|^2)$, with
\[
 f(u)=R e^{-u/2}\{I_0(u/2)+I_1(u/2)\}
 =\frac{2R}{\pi}\int_0^1e^{-us}\sqrt{\frac{1-s}{s}}\,ds.
\]
Differentiating the Gaussian convolution (a bounded function convolved
with a smooth Gaussian density) gives the score identities
\begin{align*}
 \mathbb E[(|Z|^2-1)k_w]&=w\kappa(|w|^2),\\
 \mathbb E[\overline Z^{,2}k_w]/\sqrt2
      &=\overline w\kappa(|w|^2)/\sqrt2,
 \qquad \kappa(u)=2f'(u)+uf''(u)=-\frac R2G(u).
\end{align*}
These follow equally by applying $\partial_w\partial_{\bar w}$ and
$\partial_w^2$ to $r(w)$, respectively. To check the last identity,
write $b(s)=(2/\pi)\sqrt{(1-s)/s}$. Integration by parts gives
\[
 \int_0^1(us^2-2s)b(s)e^{-us}\,ds
 =\int_0^1s^2b'(s)e^{-us}\,ds
 =-\frac1\pi\int_0^1e^{-us}\sqrt{\frac{s}{1-s}}\,ds.
\]
The boundary term $s^2b(s)e^{-us}$ vanishes at both endpoints.
Thus the visible $A/B$ projection of $k_v$ is exactly
$(a(v),\overline{a(v)}/\sqrt2)$, while its $Y,W,\nu$ projections vanish.
In particular its transverse component is zero, for \emph{every}
remote $v$, with no hypothesis on its law.

\subsubsection{Uniform inverse-distance bounds [P,C]}
If a finite positive planar measure has mass $m$ and density at most
$M/\pi$, then, at every centre $w$,
\begin{equation}\label{eq:sector-density-bound}
 \int\frac{d\mu(z)}{|z-w|}\le2\sqrt{mM}.
\end{equation}
Indeed the layer-cake integral is at most
$\int_0^\infty\min(m,M/s^2)\,ds=2\sqrt{mM}$.

First, the two orthonormal $A/B$ polynomials are
$|Z|^2-1$ and $Z^2/\sqrt2$. Their squared Christoffel sum is
$K_{AB}(s)=(s-1)^2+s^2/2$.
The cubic exponential Taylor polynomial exceeds it, since
\[
 (1+s+s^2/2+s^3/6)-K_{AB}(s)
       =s\{(s-3)^2+9\}/6\ge0.
\]
Hence $K_{AB}(s)e^{-s}\le1$. Every linear combination $P_{AB}$
therefore has density $|P_{AB}(z)|^2e^{-|z|^2}/\pi$ bounded by
$\|P_{AB}\|_2^2/\pi$, and
\[
 \mathbb E\frac{|P_{AB}(Z)|^2}{|Z+w|}\le2\|P_{AB}\|_2^2.
\]

The even-sector polynomials $Z,\overline Z,H(Z)$ are likewise
orthonormal. Their squared sum is
$K_{\rm ev}(s)=s^3/2-2s^2+4s$. Again
$K_{\rm ev}(s)e^{-s}\le1$ for all $s\ge0$.
For a short exact proof on $[0,2]$, multiply the difference between the
fourth exponential Taylor polynomial and $K_{\rm ev}$ by $24$:
\[
 Q(s)=s^4-8s^3+60s^2-72s+24.
\]
Its degree-four Bernstein coefficients are positive on each of the eight
intervals $[j/4,(j+1)/4]$, $0\le j<8$; their least coefficient is
$39/128$. The standard-library replay computes these rational
coefficients. For $s\ge2$, set $x=s-2\ge0$ and observe
\[
 -2(K_{\rm ev}'-K_{\rm ev})
 =x^3-x^2+4=(x-2/3)^2(x+1/3)+104/27>0.
\]
Thus the density decreases there and is at most $4e^{-2}<1$.
Equation~\eqref{eq:sector-density-bound} gives the same inverse-distance
constant $2$ for every even-sector combination $P_{\rm ev}$.

The transverse tests have a better constant. For a complex coefficient $a$,
put
\[
 T(a)(Z)=\frac{(|Z|^2-1)a-Z^2\bar a}{\sqrt3}.
\]
This is a \emph{real}-linear isometry: $\|T(a)\|_2^2=|a|^2$.
For $s\le1$, its density is at most $|a|^2/(3\pi)$; for $s\ge1$,
it is at most $|a|^2(2s-1)^2e^{-s}/(3\pi)$.
The latter radial expression is largest at $s=5/2$, and
\[
 \frac{16e^{-5/2}}3<\frac49,
 \qquad e^{5/2}>\sum_{j=0}^{6}(5/2)^j/j!>12.
\]
Consequently
\begin{equation}\label{eq:sector-transverse-cost}
 \mathbb E\frac{|T(a)(Z)|^2}{|Z+w|}\le\frac43|a|^2
 \qquad\hbox{uniformly in }w,a.
\end{equation}
All these inequalities can be conditioned on the full remote Gaussian
space. The coefficients and the shift may depend on exactly the same
remote variables.

\subsubsection{Support completion proves the theorems [P]}
For any full cubic test $v+U$, with remote $v$ and visible-sector $U$,
original boundedness gives
\[
 t+p\operatorname{Re}\langle v+U,F\rangle
 \le\tfrac12\mathbb E\{|Z+v+U|+|Z-v-U|\}.
\]
The global norm majorant
\[
 |z+u|\le|z|+\operatorname{Re}(\bar zu)/|z|+|u|^2/(2|z|)
\]
therefore bounds the right side by
\[
 \mathbb E\rho(|v|^2)+\operatorname{Re}\langle U,k_v\rangle
 +\frac14\mathbb E\left[|U|^2
       \left\{\frac1{|Z+v|}+\frac1{|Z-v|}\right\}\right].
\]
The singular points have Gaussian measure zero and the uniform estimates
above prove integrability. There is no clipping parameter or missing
small-ball error.

Split $U$ into its $A/B$ part and its even-sector part. The former is
even in $Z$, the latter odd, and the sum of the two denominators is even.
The quadratic cross term thus cancels \emph{exactly}, conditional on
all remote variables. The preceding bounds make the total quadratic
term at most $\|U_{AB}\|^2+\|U_{\rm ev}\|^2$.
Completing squares over all admissible remote coefficient fields gives
\[
 \|p r_A-a(v)\|^2+
 \|p r_B-\overline{a(v)}/\sqrt2\|^2+p^2E_{\rm ev}
 \le4\mathcal D_t(v;q).
\]
The real orthogonal coefficient rotation gives
\eqref{eq:sector-joint-profile}.

For the sharper content-free estimate, restrict $U_{AB}=T(a)$ with
$a\in H_{1,0}(X)$. Its real pairing with $k_v$ vanishes by the two
score identities, while its pairing with $F$ is
$\operatorname{Re}\langle a,r_\perp\rangle$.
The quadratic terms now cost at most
$(2/3)\|a\|^2+\|U_{\rm ev}\|^2$. Completing the two squares gives
\[
 \mathcal D_t(v;q)\ge\frac38p^2\|r_\perp\|^2
                      +\frac14p^2E_{\rm ev},
\]
which is \eqref{eq:sector-effective-moment}.
The choice $v=\lambda q$ and the optimum tangent multiplier
$\lambda=p/(2B_0)$ prove \eqref{eq:sector-energy-cut}.

\subsubsection{Diagnostic and remaining obligation [F,?]}
The accompanying diagnostic uses the twelve stored three-coordinate
all-sector candidates from the response-dual experiment, with their fresh
independent grade-three Gaussian buffer. It evaluates both new charges
on two independent Sobol scramblings and averages the buffer analytically.
This is a limit-model diagnostic, not a representation theorem for
arbitrary finite cubics.

At the stored row $t=0.845$, the near-Gaussian cases with raw payoffs about
$1.373952$ incur only $10^{-9}$--$10^{-8}$ of new support charge. Cases
with larger couplings incur charges up to about $1.5\cdot10^{-4}$, but
those cases already have negative scalar support gap and lower raw payoff.
Thus this sample does \emph{not} demonstrate a new endpoint gain or a
newly excluded binding candidate. The charges must not be read as changes
in the objective, and these floating values are not certificates.

What is proved is a universal replacement for leaving every shared
sector unrestricted after imposing a remote-law cut. A usable next
optimization must retain the effective first-moment constraint, the
joint parallel profile, and both existing negative squares, on the
same original-response configuration. It must then bound all high-payoff
tensors and provide a complete row cover. Those obligations remain open. \emph{Update:} the joint constraints and a compatible Ward penalty are now proved
(\S\ref{sec:joint-response-budget}); the remaining task is a universal cover of the remaining response configurations and a complete same-$\varepsilon$ row envelope. A complete
\emph{regional} row is now available (\S\ref{sec:canonical-response-cap}); a universal row is not.

\subsection{A joint support budget for the original responses}
\label{sec:joint-response-budget}
\paragraph{Provenance (IDE note, 21 October).} Lower-bound GPT joint-response-budget package (\code{conjK-joint-response-budget-20261021.zip}, archived in
\pathref{gpt-reviews/joint-response-budget-20261021/}). IDE replays \code{python -I verify\_all.py} and \code{--fine} (pinned \code{python-flint 0.9.0}, unrelated directory):
ALL CLAIMS PASS (1.5\,s and 2.6\,s), including 72 exact Ward/support audits and 18 three-branch envelope audits. IDE check \st{P} of the combination step: the transverse
test $T(a)$ is real-orthogonal to $\operatorname{ph}(Z+w)$ for every shift $w$ by the score identities, the even linear term cancels in the two-sign average by
$Z\mapsto-Z$, the inverse-distance bounds $\frac43$ and $2$ are uniform in $w$, and optimising the two test amplitudes gives $p^2/(4\cdot\frac23)=\frac{3p^2}8$ and $\frac{p^2}4$.
The uniformity in $w$ is what allows the support budgets to be added.
\emph{Scope:} the Ward-combined formula (\S\ref{sec:joint-response-budget}, last theorem) holds only at unconstrained norm maximizers in a fixed finite dimension; no
approximate-stationarity variant is claimed. The conditional $1/20000$ penalty is not a norm-cap drop and must not be added to the optimised four-block response dual.
The new example constrains only the prescribed projections onto $H$ and $J$; its fourth moment $5.93$ is a property of $q_0$ itself.


\paragraph{Scope and status.}
The results below concern the original bounded responses, not arbitrary
cubics in raw Conjecture K. They hold in every finite dimension, with
arbitrary dependence among all remote fields and without a Gaussian buffer.
[P] denotes an analytic proof and [C] a directed or exact calculation.
The global cap and the row cover remain open; no new Grothendieck endpoint
is claimed. The last theorem uses stationarity only at an unconstrained
norm maximizer in a fixed finite dimension.

Let $Z$ be the visible standard proper complex Gaussian, independent of
the remote Gaussian vector $X$. Set
\[
 h=(f_++f_-)/2,\quad k=(f_+-f_-)/2,\quad |f_\pm|\le1,
 \quad Ph=tZ,\quad F=\varepsilon Qk,\quad p=1/\varepsilon.
\]
Write
\[
 F=q+(|Z|^2-1)r_A+Z^2r_B/\sqrt2+ZY+\bar ZW+\nu H(Z),
 \qquad H(Z)=Z(|Z|^2-2)/\sqrt2,
\]
with the usual remote bidegrees. Define
\[
 r_\perp=(r_A-\sqrt2\bar r_B)/\sqrt3,\quad
 r_\parallel=(\sqrt2r_A+\bar r_B)/\sqrt3,\quad
 E_{\rm ev}=\|Y\|^2+\|W\|^2+|\nu|^2.
\]
All inner products in optimizations are real parts of the complex
$L^2$ inner product. Put $R=\sqrt\pi/2$, $\bar k=\mathbb E[k\mid X]$,
$\bar h=\mathbb E[h\mid X]$, and
\[
 \rho(u)=\mathbb E|Z+\sqrt u|,\qquad
 r(v)=2v\rho'(|v|^2),\qquad
 D(v)=\mathbb E\rho(|v|^2)-t-p\operatorname{Re}\langle v,q\rangle
 \quad(v\in H_{2,1}(X)).
\]

\subsubsection{The costs that really can be added [P]}
For every remote cubic test $v$,
\begin{equation}\label{eq:jrb-main}
 \boxed{\quad
 \frac{\|\bar k-r(v)\|^2}{\sqrt\pi}
 +\frac{3p^2}{8}\|r_\perp\|^2+\frac{p^2}{4}E_{\rm ev}
 \le D(v).\quad}
\end{equation}
There is also the complementary inequality
\begin{equation}\label{eq:jrb-h}
 \boxed{\quad
 \frac{\|\bar h\|^2+\|\bar k-r(v)\|^2}{\sqrt\pi}
 +\frac{3p^2}{8}\|r_\perp\|^2\le D(v).\quad}
\end{equation}
These strengthen the previously separate projection and sector constraints.
They do not authorize adding an arbitrary parallel-profile penalty or
adding the two displayed left sides against one copy of $D$.

\paragraph{Why the addition is valid.}
For a fixed remote shift $w$, let
\[
 k_w(Z)=\tfrac12\{\operatorname{ph}(Z+w)-\operatorname{ph}(Z-w)\},
 \qquad T(a)=\{(|Z|^2-1)a-Z^2\bar a\}/\sqrt3.
\]
The Gaussian score identities give, with a real function $\kappa$,
\[
 \mathbb E[(|Z|^2-1)\operatorname{ph}(Z+w)]=w\kappa(|w|^2),\qquad
 \mathbb E[\bar Z^2\operatorname{ph}(Z+w)]=\bar w\kappa(|w|^2).
\]
Consequently
\begin{equation}\label{eq:jrb-annihilator}
 \operatorname{Re}\mathbb E[\overline{T(a)}\operatorname{ph}(Z+w)]=0
 \qquad\hbox{for every complex }a,w.
\end{equation}
This is a real orthogonality statement; the complex pairing need not vanish.
The even-sector linear term against $k_w$ vanishes by $Z\mapsto-Z$.
The already established, uniform inverse-distance bounds are
\[
 \mathbb E\frac{|T(a)|^2}{|Z+w|}\le\frac43|a|^2,\qquad
 \mathbb E\frac{|E(Z)|^2}{|Z+w|}\le2\|E\|_{L^2(Z)}^2,
 \quad E\in\operatorname{span}\{Z,\bar Z,H(Z)\}.
\]
Their density-polynomial certificates are replayed in this package.
The norm majorant
\[
 |z+u|\le |z|+\operatorname{Re}(\bar zu)/|z|+|u|^2/(2|z|)
\]
and parity cancel the mixed transverse/even quadratic term in the
two-sign average. Completing squares therefore gives, for \emph{every}
remote $w\in L^2$, not just cubic $w$,
\begin{equation}\label{eq:jrb-shift-stable}
 t+C+\operatorname{Re}\langle w,\bar k\rangle
 \le\mathcal F(w):=\mathbb E\rho(|w|^2),\qquad
 C=\frac{3p^2}{8}\|r_\perp\|^2+\frac{p^2}{4}E_{\rm ev}.
\end{equation}
This uniformity in $w$ is the extra fact needed to combine the budgets.
There need not be an actual response with first moment $(t+C)Z$.

The real gradient of $\mathcal F$ is $r(w)$ and its Hessian is bounded
above by $R I$: $2\rho'\le R$ and $\rho''\le0$. For $w=v+u$,
\[
 D(v)-C\ge \operatorname{Re}\langle u,\bar k-r(v)\rangle
                      -\frac R2\|u\|^2.
\]
Optimizing over every remote $u$ proves \eqref{eq:jrb-main}.
For \eqref{eq:jrb-h}, use
\[
 \tfrac12\mathbb E\{\rho(|u+w|^2)+\rho(|u-w|^2)\}
 \ge t+\operatorname{Re}\langle u,\bar h\rangle
       +\operatorname{Re}\langle w,\bar k\rangle
       +3p^2\|r_\perp\|^2/8.
\]
Equation~\eqref{eq:jrb-annihilator} holds separately at the two shifts.
The product-space Hessian is again at most $R I$ and its gradient at
$(0,v)$ is $(0,r(v))$. The same completion proves the result. We do not
include the even-sector charge in this argument: the common remote shift
$u$ would destroy the parity used to eliminate its linear term.

\subsubsection{Including the parallel profile [P]}
One can also couple the nonconstant parallel profile to the projection
defect. Put
\[
 G(u)=\frac2\pi\int_0^1e^{-us}\sqrt{\frac{s}{1-s}}\,ds,\qquad
 a(v)=-\frac R2P_X\{vG(|v|^2)\},
\]
where $P_X$ is remote first-chaos projection. The map $a$ is globally
$R/2$-Lipschitz on complex $L^2$, viewed as a real Hilbert space.
Indeed $G$ is a probability mixture of $e^{-su}$, and the real derivative
eigenvalues of $w\mapsto we^{-s|w|^2}$ are $e^{-x}$ and
$(1-2x)e^{-x}$, both of absolute value at most one. Projection is a
contraction. No high derivative of the nonlinear visible profile is used.

Let
\[
 M=\|\bar k-r(v)\|,\quad
 B=\|p r_\parallel-\sqrt{3/2}\,a(v)\|,\quad
 L_0=R\sqrt{3/8},\quad K_0=R-L_0^2/2>0.
\]
Define the nondecreasing function of $M\ge0$
\[
 \Phi(M,B)=\sup_{s\ge0}
 \left\{Ms-\frac R2s^2+\frac14(B-L_0s)_+^2\right\}.
\]
It has the explicit form
\[
 \Phi(M,B)=\begin{cases}
 B^2/4,&M\le L_0B/2,\\
 B^2/4+(M-L_0B/2)^2/(2K_0),&L_0B/2<M<RB/L_0,\\
 M^2/(2R),&M\ge RB/L_0.
 \end{cases}
\]
For $B=0$ use the last branch. Then
\begin{equation}\label{eq:jrb-parallel}
 \boxed{\quad
 \Phi(M,B)+\frac{p^2}{4}(\|r_\perp\|^2+E_{\rm ev})\le D(v).
 \quad}
\end{equation}
To prove this, apply the existing full profile constraint at $v+u$,
the $R$-smooth majorant for $\mathcal F$, and
$\|a(v+u)-a(v)\|\le R\|u\|/2$. Choose $u$ in the direction of
$\bar k-r(v)$ and then optimize its length. The three branches follow
by maximizing the two quadratic pieces. In particular any certified
lower bound on $M$ can be substituted in $\Phi$. This substitution would
not generally be valid for an arbitrary, nonoptimized trial expression.
Equation~\eqref{eq:jrb-parallel} and \eqref{eq:jrb-main} are both necessary
constraints; their left sides must not be added.

\subsubsection{A joint lower bound for the two retained squares [P]}
Choose any finite remote first-chaos real subspace $\mathcal A$ and
remote cubic subspace $\mathcal B$. Set
\[
 a=\Pi_{\mathcal A}r(v),\quad
 b^2=\|\Pi_{\mathcal B}(pq-r(v))\|^2,\quad
 \tau=\sqrt\pi\{D(v)-C\}-b^2.
\]
Every feasible response has $\tau\ge0$ and
\begin{equation}\label{eq:jrb-ball}
 \|\Pi_{\mathcal A}Pk-a\|^2\le\tau,\qquad
 \|Pk\|^2\ge(\|a\|-\sqrt\tau)_+^2.
\end{equation}
The subspaces are orthogonal, so this is Bessel's inequality applied to
\eqref{eq:jrb-main}. It remains valid with uncontrolled components in
every other remote or visible direction.

At an unconstrained norm maximizer, the established visible--remote
rotation identity says
\[
 d=-\varepsilon(t x+J^*y),\quad
 x=\Pi_{\mathcal A}Pk,\quad y=Qh+B_{\rm vis}/\varepsilon,\quad
 d=\Pi_{\mathcal A}(2\beta r_A+\sqrt2\beta\bar r_B),
\]
where $B_{\rm vis}=\beta Z(|Z|^2-2)$ and $J$ is the rotation derivative
of $F$ on $\mathcal A$. Let $\Omega>0$ satisfy $\|J\|^2\le\Omega$;
$\Omega=9\|F\|^2$ is always available when $F\ne0$. Then
\begin{equation}
 \boxed{\begin{aligned}\|Pk\|^2+\varepsilon\|y\|^2\ge{}&
 \frac{\|d\|^2}{\varepsilon^2t^2+\varepsilon\Omega}\\
 &+\left(1+\frac{\varepsilon t^2}{\Omega}\right)
 \left(\left\|a+\frac{td}{\Omega+\varepsilon t^2}\right\|
                    -\sqrt\tau\right)_+^2.
 \end{aligned}}
 \label{eq:jrb-ward}
\end{equation}
This is a compatible increment to the rotation bound, not a sum of two
separately optimized penalty bounds. Indeed the left side is at least
\[
 \|x\|^2+\frac{\varepsilon}{\Omega}\|d/\varepsilon+tx\|^2
 =\frac{\|d\|^2}{\varepsilon^2t^2+\varepsilon\Omega}
 +\left(1+\frac{\varepsilon t^2}{\Omega}\right)
       \left\|x+\frac{td}{\Omega+\varepsilon t^2}\right\|^2.
\]
Minimization over the ball \eqref{eq:jrb-ball} gives \eqref{eq:jrb-ward}.
The result may be subtracted from the raw row payoff at norm maximizers.
It is not an extra subtraction from a dual which already optimizes these
same squares. A fixed-$t$ constrained row optimizer need not obey the
rotation identity.

\subsubsection{A fresh certified coupled example [P,C]}
Take independent remote standard complex Gaussians $X,Y_j,W_j$, $1\le j\le4$,
and put
\[
 H=X(|X|^2-2)/\sqrt2,\quad V=\tfrac12\sum_{j=1}^4Y_j\bar W_j,
 \quad J=XV,\quad e=38739/10^6,\quad q_0=\sqrt{e/10}(H+3J).
\]
$H,J$ are orthonormal and $\mathbb E|q_0|^4/e^2=593/100$.
Fix $t=1057/1250$, $\varepsilon=1891/2500$, $\lambda=42/25$,
and $v=\lambda q_0$. Fresh directed integration of the same positive
Gaussian kernels, including both infinite tails, proves
\[
 D(v)\simeq2.41869447\cdot10^{-5},\qquad
 b^2=\|\Pi_{\operatorname{span}\{H,J\}}(pq_0-r(v))\|^2
       \simeq4.00139386\cdot10^{-5},
\]
and, rigorously,
\begin{equation}\label{eq:jrb-budget-number}
 0<D(v)-b^2/\sqrt\pi<1612/10^9.
\end{equation}
Every scalar support cut for $q_0$ passes with margin $>1/50000$.
Nevertheless any original response whose remote coefficients along
$H,J$ equal those of $q_0$ must obey
\[
 \frac{3p^2}{8}\|r_\perp\|^2+\frac{p^2}{4}E_{\rm ev}
       <1612/10^9,
 \qquad \|r_\perp\|<0.001569,\quad E_{\rm ev}<0.000003689.
\]
All orthogonal remote components are unrestricted in this assertion.
In particular $\|r_\perp\|\ge1/500$ is impossible, with a support-budget
violation exceeding $10^{-6}$. At $1/500$, the displayed projection test
and displayed sector test each pass separately. This does not assert
that the configuration passes every other previously available test.

For the same prescribed two remote coefficients, if feasible,
\[
 \|r_\perp\|\ge1/1000\ \Longrightarrow\ \|Pk\|^2>1/60000,
 \qquad
 \|r_\perp\|\ge3/2000\ \Longrightarrow\ \|Pk\|^2>1/42000.
\]
The implication includes infeasibility when $\tau<0$; at the endpoints
$\tau>0$ and the checker completes the squares directly.

Finally take exactly $q=q_0$, $r_\parallel=-X/1000$,
$r_\perp=X/1000$, and $Y=W=\nu=0$. Thus
$r_A=(\sqrt2r_\parallel+r_\perp)/\sqrt3$ and
$\bar r_B=(r_\parallel-\sqrt2r_\perp)/\sqrt3$.
If this configuration occurs at an unconstrained norm maximizer, then
\[
 \boxed{\ \|Pk\|^2+\varepsilon\|Qh+B_{\rm vis}/\varepsilon\|^2
            >1/20000.\ }
\]
The computed joint lower bound exceeds $0.00005121$;
the standalone rotation bound with the same $\Omega=9(e+2\cdot10^{-6})$
is below $0.0000002$. This is a conditional necessary penalty, not a
construction of a maximizer, nor a norm-cap improvement of $1/20000$.

\paragraph{Next numerical obligation.}
Use \eqref{eq:jrb-main}, \eqref{eq:jrb-parallel}, and the full response
dual as constraints on the same tensor. At a norm-maximizer relaxation,
\eqref{eq:jrb-ward} is an alternative certified lower estimate of its
existing squares. A universal cap still requires a dimension-free
reduction or exhaustive cover of feasible high-payoff configurations,
and a complete row cover at one common $\varepsilon$.

